\section*{Hoofdstuk \ref{ch:g9:solids} --- Lichamen en inhouden}

\begin{solution}{exo:g9:solids:1}
\omterm{def:g5:solids:def}{Balk}: $V = 4 \times 5 \times 12 = 240$.

\omterm{def:g7:areas:prism}{Cilinder}: $V = \pi r^2 h = \pi \times 9 \times 7 = 63\pi \approx 198$.
\end{solution}

\begin{solution}{exo:g9:solids:2}
\omterm{def:g8:solids:def}{Kegel}: $V = \frac13 \pi r^2 h = \frac13 \pi \times 25 \times 12 = 100\pi
\approx 314$.

\omterm{def:g8:solids:def}{Piramide}: \omterm{def:g6:measure:area}{oppervlakte} van het grondvlak $B = 8 \times 3 = 24$, dus
$V = \frac13 \times 24 \times 10 = 80$.
\end{solution}

\begin{solution}{exo:g9:solids:3}
\omterm{def:g2:measure:units}{Inhoud}: $V = \frac43 \pi \times 6^3 = \frac43 \pi \times 216 = 288\pi$.

\omterm{def:g6:measure:area}{Oppervlakte} van het boloppervlak: $4\pi \times 6^2 = 144\pi$.
\end{solution}

\begin{solution}{exo:g9:solids:4}
\omterm{def:g3:shapes:shapes}{Radius} van de \omterm{prop:g9:solids:sections}{doorsnede}: $\sqrt{r^2 - d^2} = \sqrt{100 - 64} = \sqrt{36} = 6$.
\end{solution}

\begin{solution}{exo:g9:solids:5}
\omterm{def:g9:thales:scaling}{Schaalfactor} van de hele \omterm{def:g8:solids:def}{kegel} naar de kleine: $k = \frac{9}{12} = \frac34$.
\omterm{def:g3:shapes:shapes}{Radius} van de \omterm{prop:g9:solids:sections}{doorsnede}: $8 \times \frac34 = 6$.

\omterm{def:g2:measure:units}{Inhoud} van de hele \omterm{def:g8:solids:def}{kegel}: $\frac13\pi \times 64 \times 12 = 256\pi$. \omterm{def:g2:measure:units}{Inhoud} van de
kleine \omterm{def:g8:solids:def}{kegel}: $256\pi \times \left(\frac34\right)^3 = 256\pi \times \frac{27}{64}
= 108\pi \approx 339$.
\end{solution}

\begin{solution}{exo:g9:solids:6}
\omterm{def:g2:measure:units}{Inhoud} van de bal: $\frac43\pi \times 2^3 = \frac{32\pi}{3}$ cm$^3$. Die \omterm{def:g2:measure:units}{inhoud}
verdeelt zich over de \omterm{def:g6:measure:area}{oppervlakte} van het grondvlak $\pi \times 3^2 = 9\pi$
cm$^2$, dus stijgt het peil met
\[
\frac{32\pi/3}{9\pi} = \frac{32}{27} \approx 1.2 \text{ cm}
\]
(ongeveer $12$ mm).
\end{solution}

\begin{solution}{exo:g9:solids:7}
De kleine vorm is de verkleining van de grote met $k = \frac36 = \frac12$, dus is
haar \omterm{def:g2:measure:units}{inhoud} $\left(\frac12\right)^3 = \frac18$ van die van de grote: het recept
vult $8$ kleine vormen.
\end{solution}

\begin{solution}{exo:g9:solids:8}
$V = \frac13 \times 230^2 \times 147 = \frac13 \times 52\,900 \times 147
= 52\,900 \times 49 = 2\,592\,100$ m$^3$ --- ongeveer $2.6$ miljoen kubieke
meter.
\end{solution}

\begin{solution}{exo:g9:solids:9}
\emph{1.} Het water vormt een \omterm{def:g8:solids:def}{kegel} geschaald met $k = \frac12$ (top naar onder,
de \omterm{def:g3:division:half}{helft} van de hoogte), dus is zijn \omterm{def:g2:measure:units}{inhoud} $\left(\frac12\right)^3 = \frac18$ van
die van de trechter: een achtste, veel minder dan de \omterm{def:g3:division:half}{helft}!

\emph{2.} Water tot hoogte $d$ vormt een \omterm{def:g8:solids:def}{kegel} geschaald met $k = \frac{d}{12}$,
met een \omterm{def:g2:measure:units}{inhoud} die $k^3$ keer die van de trechter is. De \omterm{def:g3:division:half}{helft} van de \omterm{def:g2:measure:units}{inhoud}
betekent $k^3 = \frac12$, dat wil zeggen $\left(\frac{d}{12}\right)^3 = \frac12$.
Dan is $\frac{d}{12} = \sqrt[3]{0.5} \approx 0.7937$, dus $d \approx 9.5$ cm: de
tweede \omterm{def:g3:division:half}{helft} van de \omterm{def:g2:measure:units}{inhoud} neemt maar de bovenste $2.5$ cm in --- het brede deel
van de \omterm{def:g8:solids:def}{kegel} bevat het meeste water.
\end{solution}

\begin{solution}{pb:g9:solids:1}
\textbf{1.} $V_{\text{cil}} = \pi r^2 \times 2r = 2\pi r^3$.

\textbf{2.} $\dfrac{V_{\text{bol}}}{V_{\text{cil}}}
= \dfrac{\frac43 \pi r^3}{2 \pi r^3} = \dfrac23$. De $\pi$ en de $r^3$ vallen weg:
de verhouding is \emph{universeel} --- waar voor een knikker en voor een planeet.
Een verband tussen vormen, niet tussen getallen: precies het soort waarheid dat het
verdient in steen gebeiteld te worden.

\textbf{3.} Bol: $4\pi r^2$. \omterm{def:g7:areas:prism}{Cilinder}: het \omterm{def:g5:solids:def}{zijvlak}
$2\pi r \times 2r = 4\pi r^2$, plus twee deksels $2 \times \pi r^2$: samen
$6\pi r^2$. Verhouding: $\frac{4\pi r^2}{6\pi r^2} = \frac23$ --- dezelfde twee
derde, nu voor \omterm{def:g6:measure:area}{oppervlakten}. (En een extraatje: de \omterm{def:g6:measure:area}{oppervlakte} van de bol is precies
gelijk aan die van het \omterm{def:g5:solids:def}{zijvlak} van de \omterm{def:g7:areas:prism}{cilinder}.)

\textbf{4.} Wat overblijft: $2\pi r^3 - \frac43\pi r^3 = \frac23 \pi r^3$. Twee
\omterm{def:g8:solids:def}{kegels} met \omterm{def:g3:shapes:shapes}{radius} $r$ en hoogte $r$: $2 \times \frac13 \pi r^2 \times r
= \frac23 \pi r^3$. Gelijk.

\textbf{5.} Bal: $\frac43 \pi \times 12^3 \approx 7\,200$ cm$^3$; doos:
$2\pi \times 12^3 \approx 10\,900$ cm$^3$. Leeg: een derde van de doos, ongeveer
$33\,\%$ --- gewaarborgd door vraag 2, wat de grootte van de bal ook is.

\textbf{6.} De \omterm{def:g3:division:half}{helft} van een bol:
$\frac12 \times \frac43 \pi \times 10^3 = \frac23 \pi \times 1\,000 \approx
2\,094$ cm$^3 \approx 2.1$ L.

\textbf{7.} \omterm{def:g8:solids:def}{Kegel}: $\frac13 \pi \times 10^2 \times 10 = \frac{1\,000\pi}{3}
\approx 1\,047$ cm$^3$. Halve bol: $\frac{2\,000\pi}{3} \approx 2\,094$ cm$^3$.
\omterm{def:g7:areas:prism}{Cilinder}: $1\,000\pi \approx 3\,142$ cm$^3$. Verhoudingen:
$\frac{1000\pi}{3} : \frac{2000\pi}{3} : \frac{3000\pi}{3} = 1 : 2 : 3$, precies.

\textbf{8.} $\frac43 \pi \times 27 = \pi \times 9 \times h$ geeft
$h = \frac43 \times 3 = 4$ cm, precies.

\textbf{9.} Radiussen: $\frac{6\,371}{1\,737} \approx 3.67$. \omterm{def:g2:measure:units}{Inhouden} schalen met
de derde \omterm{def:g9:fractions:power}{macht} (\cref{thm:g9:solids:scaling}): $3.67^3 \approx 49$. Er passen
ongeveer vijftig Manen in de Aarde.

\textbf{10.} \omterm{def:g6:measure:area}{Oppervlakten} schalen met het kwadraat: $3.67^2 \approx 13.5$. De
\omterm{def:g3:shapes:shapes}{radius} van een ballon verdubbelen vermenigvuldigt zijn \omterm{def:g2:measure:units}{inhoud} met $2^3 = 8$, maar
zijn oppervlak met slechts $2^2 = 4$: als dingen groeien, loopt de \omterm{def:g2:measure:units}{inhoud} (gewicht,
vulling, geproduceerde warmte) het oppervlak (huid, materiaal, verloren warmte)
voorbij --- de kwadraat-kubuswet.

\textbf{11.} Gewicht: $\times 10^3 = 1\,000$. \omterm{prop:g9:solids:sections}{Doorsnede} van het been:
$\times 10^2 = 100$. Druk --- gewicht per \omterm{def:g6:measure:area}{oppervlakte} been:
$\times \frac{1000}{100} = 10$. Beenderen die voor menselijke druk gebouwd zijn,
krijgen er tien keer zoveel: het skelet van de reus knapt onder zijn eigen gewicht.
Reuzen zijn meetkundig onmogelijk; grote dieren hebben onevenredig dikke beenderen
nodig (vergelijk de benen van een olifant met die van een gazelle).

\textbf{12.} Gewicht: $\div 100^3 = 10^6$. Kracht (\omterm{prop:g9:solids:sections}{doorsnede} van de spieren):
$\div 100^2 = 10^4$. Kracht ten opzichte van gewicht:
$\times \frac{10^6}{10^4} = 100$. De mier die vijftig keer haar gewicht tilt is
geen superatleet --- ze is alleen maar klein; een mens die tot mierengrootte
krimpt, zou hetzelfde kunnen.

\textbf{13.} $\dfrac{S}{V} = \dfrac{4\pi r^2}{\frac43\pi r^3} = \dfrac{3}{r}$:
voor $r = 1$ is de verhouding $3$, voor $r = 2$ is ze $1.5$ --- de grootte
halveren verdubbelt het beschikbare oppervlak per eenheid warmteproducerende
\omterm{def:g2:measure:units}{inhoud}. Kleine lichamen bloeden warmte weg: de muis moet onophoudelijk eten, de
beer kan een winterslaap houden, en de baby --- kleine bol, grote $\frac Sr$ ---
heeft de muts nodig.

\textbf{14.} $\left(\frac54\right)^3 = \frac{125}{64} \approx 1.95$: de grote
sinaasappel bevat bijna \emph{twee keer} zoveel vrucht voor dezelfde prijs. Koop
\omterm{def:g3:shapes:shapes}{radius}.

\textbf{15.} Het beeldhouwwerk zegt: de bol is $\frac23$ van de \omterm{def:g7:areas:prism}{cilinder} in \omterm{def:g2:measure:units}{inhoud},
\emph{en} $\frac23$ ervan in totale \omterm{def:g6:measure:area}{oppervlakte} --- twee exacte, universele
verhoudingen die het rondste \omterm{def:g5:solids:def}{lichaam} met het eenvoudigste verbinden. Het antwoord
aan Cicero: omdat een stelling het enige monument is dat legers noch eeuwen
afslijten --- de machines van Archimedes verbrandden met Syracuse, maar de twee
derde is nog steeds precies twee derde.
\end{solution}
